\section[Canonical kappa-r normalization]{$\kapr$ normalization}
\label{sec:kappa-normalization}

The normalization constant \(\kapr\) is the scalar that makes the cascade
leading-term form use the same arithmetic factor as the standard Strong BSD
leading-coefficient formula.  The point of this section is only to identify
that scalar after the cascade convention \(c_E=\Tam(E)\) is fixed.  It is an
equivalence of two ways of writing the same conjectural identity, not a
derivation of BSD.

\begin{theorem}[Canonical $\kapr$ normalization equivalence]\label{thm:apparatus:kappa-r-normalization-equivalence}
Fix the cascade convention
\[
  c_E:=\Tam(E)=\prod_{\ell} c_{\ell}(E).
\]
Then the cascade leading-term form
\[
  \frac{L^{(r)}(E,1)}{r!}
  =
  \kapr(E)\,\Omega_E\,\Reg_r(E)\,c_E
\]
agrees with the standard Strong BSD leading-coefficient form
\[
  \frac{L^{(r)}(E,1)}{r!}
  =
  \Omega_E\,\Reg_r(E)\,
  \frac{\#\Sha(E/\Q)\,\Tam(E)}{\#E(\Q)_{\tors}^{\,2}}
\]
if and only if
\[
  \kapr(E)=
  \frac{\#\Sha(E/\Q)}{\#E(\Q)_{\tors}^{\,2}}.
\]
Under this canonical normalization, the cascade form is equivalent to the
standard Strong BSD identity as a restatement in cascade target coordinates.
It does not prove the identity.\footnote{\raggedright Verified in Lean as
\texttt{SixBirdsBSD.\allowbreak Apparatus.\allowbreak KappaNormalization.\allowbreak kappaRNormalizationEquivalence};
see App~\ref{app:formalization}.}
\end{theorem}

\begin{proof}
After the convention \(c_E=\Tam(E)\), the right-hand side of the cascade
form is
\[
  \kapr(E)\,\Omega_E\,\Reg_r(E)\,\Tam(E).
\]
The right-hand side of the standard Strong BSD leading-coefficient formula
is
\[
  \Omega_E\,\Reg_r(E)\,
  \frac{\#\Sha(E/\Q)\,\Tam(E)}{\#E(\Q)_{\tors}^{\,2}}
  =
  \frac{\#\Sha(E/\Q)}{\#E(\Q)_{\tors}^{\,2}}\,
  \Omega_E\,\Reg_r(E)\,\Tam(E).
\]
Thus the two formulas have the same left-hand side and differ only in the
coefficient multiplying the common factor
\[
  \Omega_E\,\Reg_r(E)\,\Tam(E).
\]
In the BSD normalization this factor is nonzero: the period is positive,
the regulator determinant is positive in positive rank and equal to \(1\)
in rank \(0\), and the Tamagawa product is a positive integer.  Therefore
the two right-hand sides are equal if and only if their scalar coefficients
are equal, namely if and only if
\[
  \kapr(E)
  =
  \frac{\#\Sha(E/\Q)}{\#E(\Q)_{\tors}^{\,2}}.
\]

If this equality of coefficients is imposed, substituting it into the
cascade form gives exactly the standard Strong BSD leading-coefficient
formula.  Conversely, if the cascade form and the standard form agree under
the fixed convention \(c_E=\Tam(E)\), cancellation of the same nonzero common
factor forces the displayed value of \(\kapr(E)\).  The proof is only this
algebraic comparison of normalizations; it supplies no proof that either
leading-coefficient identity is true for \(E\).
\end{proof}

The conditional Strong-BSD closure of \citet{TsiokosBSDClosure2026}
consumes this normalization.  The role of the present section is only to
supply the normalization equivalence used there.
